% !TeX root = ../../protocolo.tex

To investigate the spectral response of a metasurface composed of AuNPs immersed in an aqueous environment, an ATR-based setup, as depicted in Fig. \ref{fgr:Setup}, was utilized. The setup incorporates a polydimethylsiloxane (PDMS) microfluidic device positioned atop the AuNP ensemble, which is anchored to a glass substrate. This configuration enables precise control over the water flow surrounding the nanoparticles.

For optical coupling, the metasurface’s substrate was index-matched to a BK7 glass prism (P) using immersion oil, ensuring an effective refractive index of. Illumination was provided by a broad-spectrum Xenon lamp (ThorLabs SLS201L), which was first collimated using a lens (L) and subsequently passed through a linear polarizer (LP), allowing the selection of either ss- or pp-polarized light.

The sample was illuminated in an ATR configuration, meaning the angle of incidence exceeded the critical angle of the BK7-water interface. Under these conditions, an evanescent wave is generated at the interface, interacting with the AuNPs of the metasurface and thereby exciting localized surface plasmon resonances (LSPRs) on each nanoparticle.

The LSPR response was detected by analyzing the reflection spectra of the light reflected from the metasurface. This reflected beam was collected by a second lens and directed into a spectrometer (S) for spectral analysis. To ensure accurate measurements, the recorded spectra were normalized against the reflection intensity obtained from a pristine glass substrate.